% LingBot-VLA 2.0 优化报告(精简结构版)
% 编译:tectonic -X compile lingbot_vla_v2_total_report.tex
\documentclass[11pt]{ctexart}
% cnreport-preamble.tex — 中文技术报告导言区(Cosmos/arXiv 单栏风格)
% 用法: \documentclass[11pt]{ctexart} 之后 % cnreport-preamble.tex — 中文技术报告导言区(Cosmos/arXiv 单栏风格)
% 用法: \documentclass[11pt]{ctexart} 之后 % cnreport-preamble.tex — 中文技术报告导言区(Cosmos/arXiv 单栏风格)
% 用法: \documentclass[11pt]{ctexart} 之后 \input{cnreport-preamble}
% 编译: tectonic(XeTeX 引擎),需要系统安装 Noto Serif/Sans CJK SC 字体
\usepackage[a4paper,top=2.7cm,bottom=2.7cm,left=2.5cm,right=2.5cm,headheight=15pt]{geometry}
\usepackage{fontspec}
% 正文:拉丁衬线 + 思源宋体;标题:拉丁无衬线 + 思源黑体
% 按文件名从 TeX 字体树加载(tectonic bundle 自带 TeX Gyre;系统字体需 Noto CJK SC)
\setmainfont{texgyretermes}[Extension=.otf,UprightFont=*-regular,BoldFont=*-bold,ItalicFont=*-italic,BoldItalicFont=*-bolditalic]
\setsansfont{texgyreheros}[Extension=.otf,UprightFont=*-regular,BoldFont=*-bold,ItalicFont=*-italic,BoldItalicFont=*-bolditalic]
\setmonofont{texgyrecursor}[Extension=.otf,UprightFont=*-regular,BoldFont=*-bold,Scale=0.88]
\setCJKmainfont{Noto Serif CJK SC}[BoldFont={Noto Serif CJK SC Bold}]
\setCJKsansfont{Noto Sans CJK SC}[BoldFont={Noto Sans CJK SC Bold}]
\setCJKmonofont{Noto Sans Mono CJK SC}
% 带圈数字 ①-⑳ 等划入 CJK 类,用 CJK 字体渲染(xeCJK 默认按拉丁处理)
\xeCJKDeclareCharClass{CJK}{"2460 -> "24FF}

\usepackage{amsmath,amssymb,bm}
\usepackage{booktabs}
\usepackage{graphicx}
\usepackage{xcolor}
\usepackage{titlesec}
\usepackage{fancyhdr}
\usepackage{enumitem}
\usepackage{caption}
\usepackage{listings}

% ---- 配色(主色可改) ----
\definecolor{accent}{HTML}{0E5FA8}     % 强调蓝(链接/节号)
\definecolor{good}{HTML}{0E7A4D}       % 正向结果
\definecolor{warn}{HTML}{B45309}       % 负向/警示
\definecolor{faint}{HTML}{6B7889}      % 弱化文字
\definecolor{codebg}{HTML}{F0F2F5}
\definecolor{rule}{HTML}{C9CFD8}

\usepackage[colorlinks=true,linkcolor=accent,citecolor=accent,urlcolor=accent]{hyperref}

% ---- 章节标题:无衬线粗体,节号着色 ----
\titleformat{\section}{\Large\bfseries\sffamily}{\textcolor{accent}{\thesection}}{0.6em}{}
\titleformat{\subsection}{\large\bfseries\sffamily}{\textcolor{accent}{\thesubsection}}{0.5em}{}
\titleformat{\subsubsection}{\normalsize\bfseries\sffamily}{\textcolor{accent}{\thesubsubsection}}{0.5em}{}
\titlespacing{\section}{0pt}{1.6em}{0.7em}
\titlespacing{\subsection}{0pt}{1.2em}{0.5em}

% ---- 页眉页脚(Cosmos 式:左标题右日期,下细线) ----
\newcommand{\runningtitle}{}  % 主文档用 \renewcommand 设置
\newcommand{\reportdate}{}
\newcommand{\reportfooter}{}
\fancypagestyle{cnreport}{%
  \fancyhf{}
  \fancyhead[L]{\small\sffamily \runningtitle}
  \fancyhead[R]{\small\sffamily \reportdate}
  \fancyfoot[L]{\small\color{faint}\sffamily \reportfooter}
  \fancyfoot[R]{\small\sffamily \thepage}
  \renewcommand{\headrulewidth}{0.4pt}
  \renewcommand{\footrulewidth}{0pt}
}
\pagestyle{cnreport}

% ---- 题注与列表 ----
\captionsetup{font=small,labelfont={bf,sf},skip=6pt}
\setlist{itemsep=2pt,topsep=4pt,parsep=0pt}

% ---- 代码 ----
\lstset{
  basicstyle=\ttfamily\small,
  backgroundcolor=\color{codebg},
  frame=single, framerule=0.3pt, rulecolor=\color{rule},
  xleftmargin=6pt, xrightmargin=6pt,
  breaklines=true, columns=fullflexible, keepspaces=true,
  aboveskip=8pt, belowskip=8pt,
}

% ---- 报告头(kicker / 标题 / 副题 / 署名 / 摘要 / 关键词) ----
% 用法: \reporthead{kicker-left}{kicker-right}{标题}{副题}{署名行}
\newcommand{\reporthead}[5]{%
  \begin{center}
  {\small\sffamily\bfseries\color{accent} #1 \hfill #2}\\[2pt]
  {\color{black}\rule{\linewidth}{1.2pt}}\\[14pt]
  {\LARGE\sffamily\bfseries #3\par}\vspace{8pt}
  {\large\color{faint} #4\par}\vspace{10pt}
  {\small\sffamily\color{faint} #5\par}\vspace{4pt}
  {\color{rule}\rule{\linewidth}{0.4pt}}
  \end{center}\vspace{6pt}
}
% 摘要块: \begin{cnabstract} ... \end{cnabstract}
\newenvironment{cnabstract}{%
  \begin{quote}\small\color{faint}\sffamily\bfseries\color{black} 摘要 \quad\normalfont\rmfamily\color{faint!80!black}%
}{\end{quote}\vspace{2pt}}
% 重点结论框: \begin{quotebox} ... \end{quotebox}
\newenvironment{quotebox}{%
  \begin{quote}\small
  {\color{rule}\rule{\linewidth}{0.8pt}}\par\vspace{2pt}%
}{\par\vspace{2pt}{\color{rule}\rule{\linewidth}{0.8pt}}\end{quote}\vspace{2pt}}

\newcommand{\keywords}[1]{\begin{quote}\small\color{faint}{\bfseries\sffamily\color{black} 关键词:\ } #1\end{quote}}

% 表内彩色标记
\newcommand{\upgood}[1]{\textcolor{good}{\bfseries #1}}
\newcommand{\downbad}[1]{\textcolor{warn}{\bfseries #1}}
\newcommand{\dimtext}[1]{\textcolor{faint}{#1}}

\setlength{\parskip}{2pt}
\linespread{1.12}
\emergencystretch=3em  % CJK 段落容忍度,消除 overfull

% 编译: tectonic(XeTeX 引擎),需要系统安装 Noto Serif/Sans CJK SC 字体
\usepackage[a4paper,top=2.7cm,bottom=2.7cm,left=2.5cm,right=2.5cm,headheight=15pt]{geometry}
\usepackage{fontspec}
% 正文:拉丁衬线 + 思源宋体;标题:拉丁无衬线 + 思源黑体
% 按文件名从 TeX 字体树加载(tectonic bundle 自带 TeX Gyre;系统字体需 Noto CJK SC)
\setmainfont{texgyretermes}[Extension=.otf,UprightFont=*-regular,BoldFont=*-bold,ItalicFont=*-italic,BoldItalicFont=*-bolditalic]
\setsansfont{texgyreheros}[Extension=.otf,UprightFont=*-regular,BoldFont=*-bold,ItalicFont=*-italic,BoldItalicFont=*-bolditalic]
\setmonofont{texgyrecursor}[Extension=.otf,UprightFont=*-regular,BoldFont=*-bold,Scale=0.88]
\setCJKmainfont{Noto Serif CJK SC}[BoldFont={Noto Serif CJK SC Bold}]
\setCJKsansfont{Noto Sans CJK SC}[BoldFont={Noto Sans CJK SC Bold}]
\setCJKmonofont{Noto Sans Mono CJK SC}
% 带圈数字 ①-⑳ 等划入 CJK 类,用 CJK 字体渲染(xeCJK 默认按拉丁处理)
\xeCJKDeclareCharClass{CJK}{"2460 -> "24FF}

\usepackage{amsmath,amssymb,bm}
\usepackage{booktabs}
\usepackage{graphicx}
\usepackage{xcolor}
\usepackage{titlesec}
\usepackage{fancyhdr}
\usepackage{enumitem}
\usepackage{caption}
\usepackage{listings}

% ---- 配色(主色可改) ----
\definecolor{accent}{HTML}{0E5FA8}     % 强调蓝(链接/节号)
\definecolor{good}{HTML}{0E7A4D}       % 正向结果
\definecolor{warn}{HTML}{B45309}       % 负向/警示
\definecolor{faint}{HTML}{6B7889}      % 弱化文字
\definecolor{codebg}{HTML}{F0F2F5}
\definecolor{rule}{HTML}{C9CFD8}

\usepackage[colorlinks=true,linkcolor=accent,citecolor=accent,urlcolor=accent]{hyperref}

% ---- 章节标题:无衬线粗体,节号着色 ----
\titleformat{\section}{\Large\bfseries\sffamily}{\textcolor{accent}{\thesection}}{0.6em}{}
\titleformat{\subsection}{\large\bfseries\sffamily}{\textcolor{accent}{\thesubsection}}{0.5em}{}
\titleformat{\subsubsection}{\normalsize\bfseries\sffamily}{\textcolor{accent}{\thesubsubsection}}{0.5em}{}
\titlespacing{\section}{0pt}{1.6em}{0.7em}
\titlespacing{\subsection}{0pt}{1.2em}{0.5em}

% ---- 页眉页脚(Cosmos 式:左标题右日期,下细线) ----
\newcommand{\runningtitle}{}  % 主文档用 \renewcommand 设置
\newcommand{\reportdate}{}
\newcommand{\reportfooter}{}
\fancypagestyle{cnreport}{%
  \fancyhf{}
  \fancyhead[L]{\small\sffamily \runningtitle}
  \fancyhead[R]{\small\sffamily \reportdate}
  \fancyfoot[L]{\small\color{faint}\sffamily \reportfooter}
  \fancyfoot[R]{\small\sffamily \thepage}
  \renewcommand{\headrulewidth}{0.4pt}
  \renewcommand{\footrulewidth}{0pt}
}
\pagestyle{cnreport}

% ---- 题注与列表 ----
\captionsetup{font=small,labelfont={bf,sf},skip=6pt}
\setlist{itemsep=2pt,topsep=4pt,parsep=0pt}

% ---- 代码 ----
\lstset{
  basicstyle=\ttfamily\small,
  backgroundcolor=\color{codebg},
  frame=single, framerule=0.3pt, rulecolor=\color{rule},
  xleftmargin=6pt, xrightmargin=6pt,
  breaklines=true, columns=fullflexible, keepspaces=true,
  aboveskip=8pt, belowskip=8pt,
}

% ---- 报告头(kicker / 标题 / 副题 / 署名 / 摘要 / 关键词) ----
% 用法: \reporthead{kicker-left}{kicker-right}{标题}{副题}{署名行}
\newcommand{\reporthead}[5]{%
  \begin{center}
  {\small\sffamily\bfseries\color{accent} #1 \hfill #2}\\[2pt]
  {\color{black}\rule{\linewidth}{1.2pt}}\\[14pt]
  {\LARGE\sffamily\bfseries #3\par}\vspace{8pt}
  {\large\color{faint} #4\par}\vspace{10pt}
  {\small\sffamily\color{faint} #5\par}\vspace{4pt}
  {\color{rule}\rule{\linewidth}{0.4pt}}
  \end{center}\vspace{6pt}
}
% 摘要块: \begin{cnabstract} ... \end{cnabstract}
\newenvironment{cnabstract}{%
  \begin{quote}\small\color{faint}\sffamily\bfseries\color{black} 摘要 \quad\normalfont\rmfamily\color{faint!80!black}%
}{\end{quote}\vspace{2pt}}
% 重点结论框: \begin{quotebox} ... \end{quotebox}
\newenvironment{quotebox}{%
  \begin{quote}\small
  {\color{rule}\rule{\linewidth}{0.8pt}}\par\vspace{2pt}%
}{\par\vspace{2pt}{\color{rule}\rule{\linewidth}{0.8pt}}\end{quote}\vspace{2pt}}

\newcommand{\keywords}[1]{\begin{quote}\small\color{faint}{\bfseries\sffamily\color{black} 关键词:\ } #1\end{quote}}

% 表内彩色标记
\newcommand{\upgood}[1]{\textcolor{good}{\bfseries #1}}
\newcommand{\downbad}[1]{\textcolor{warn}{\bfseries #1}}
\newcommand{\dimtext}[1]{\textcolor{faint}{#1}}

\setlength{\parskip}{2pt}
\linespread{1.12}
\emergencystretch=3em  % CJK 段落容忍度,消除 overfull

% 编译: tectonic(XeTeX 引擎),需要系统安装 Noto Serif/Sans CJK SC 字体
\usepackage[a4paper,top=2.7cm,bottom=2.7cm,left=2.5cm,right=2.5cm,headheight=15pt]{geometry}
\usepackage{fontspec}
% 正文:拉丁衬线 + 思源宋体;标题:拉丁无衬线 + 思源黑体
% 按文件名从 TeX 字体树加载(tectonic bundle 自带 TeX Gyre;系统字体需 Noto CJK SC)
\setmainfont{texgyretermes}[Extension=.otf,UprightFont=*-regular,BoldFont=*-bold,ItalicFont=*-italic,BoldItalicFont=*-bolditalic]
\setsansfont{texgyreheros}[Extension=.otf,UprightFont=*-regular,BoldFont=*-bold,ItalicFont=*-italic,BoldItalicFont=*-bolditalic]
\setmonofont{texgyrecursor}[Extension=.otf,UprightFont=*-regular,BoldFont=*-bold,Scale=0.88]
\setCJKmainfont{Noto Serif CJK SC}[BoldFont={Noto Serif CJK SC Bold}]
\setCJKsansfont{Noto Sans CJK SC}[BoldFont={Noto Sans CJK SC Bold}]
\setCJKmonofont{Noto Sans Mono CJK SC}
% 带圈数字 ①-⑳ 等划入 CJK 类,用 CJK 字体渲染(xeCJK 默认按拉丁处理)
\xeCJKDeclareCharClass{CJK}{"2460 -> "24FF}

\usepackage{amsmath,amssymb,bm}
\usepackage{booktabs}
\usepackage{graphicx}
\usepackage{xcolor}
\usepackage{titlesec}
\usepackage{fancyhdr}
\usepackage{enumitem}
\usepackage{caption}
\usepackage{listings}

% ---- 配色(主色可改) ----
\definecolor{accent}{HTML}{0E5FA8}     % 强调蓝(链接/节号)
\definecolor{good}{HTML}{0E7A4D}       % 正向结果
\definecolor{warn}{HTML}{B45309}       % 负向/警示
\definecolor{faint}{HTML}{6B7889}      % 弱化文字
\definecolor{codebg}{HTML}{F0F2F5}
\definecolor{rule}{HTML}{C9CFD8}

\usepackage[colorlinks=true,linkcolor=accent,citecolor=accent,urlcolor=accent]{hyperref}

% ---- 章节标题:无衬线粗体,节号着色 ----
\titleformat{\section}{\Large\bfseries\sffamily}{\textcolor{accent}{\thesection}}{0.6em}{}
\titleformat{\subsection}{\large\bfseries\sffamily}{\textcolor{accent}{\thesubsection}}{0.5em}{}
\titleformat{\subsubsection}{\normalsize\bfseries\sffamily}{\textcolor{accent}{\thesubsubsection}}{0.5em}{}
\titlespacing{\section}{0pt}{1.6em}{0.7em}
\titlespacing{\subsection}{0pt}{1.2em}{0.5em}

% ---- 页眉页脚(Cosmos 式:左标题右日期,下细线) ----
\newcommand{\runningtitle}{}  % 主文档用 \renewcommand 设置
\newcommand{\reportdate}{}
\newcommand{\reportfooter}{}
\fancypagestyle{cnreport}{%
  \fancyhf{}
  \fancyhead[L]{\small\sffamily \runningtitle}
  \fancyhead[R]{\small\sffamily \reportdate}
  \fancyfoot[L]{\small\color{faint}\sffamily \reportfooter}
  \fancyfoot[R]{\small\sffamily \thepage}
  \renewcommand{\headrulewidth}{0.4pt}
  \renewcommand{\footrulewidth}{0pt}
}
\pagestyle{cnreport}

% ---- 题注与列表 ----
\captionsetup{font=small,labelfont={bf,sf},skip=6pt}
\setlist{itemsep=2pt,topsep=4pt,parsep=0pt}

% ---- 代码 ----
\lstset{
  basicstyle=\ttfamily\small,
  backgroundcolor=\color{codebg},
  frame=single, framerule=0.3pt, rulecolor=\color{rule},
  xleftmargin=6pt, xrightmargin=6pt,
  breaklines=true, columns=fullflexible, keepspaces=true,
  aboveskip=8pt, belowskip=8pt,
}

% ---- 报告头(kicker / 标题 / 副题 / 署名 / 摘要 / 关键词) ----
% 用法: \reporthead{kicker-left}{kicker-right}{标题}{副题}{署名行}
\newcommand{\reporthead}[5]{%
  \begin{center}
  {\small\sffamily\bfseries\color{accent} #1 \hfill #2}\\[2pt]
  {\color{black}\rule{\linewidth}{1.2pt}}\\[14pt]
  {\LARGE\sffamily\bfseries #3\par}\vspace{8pt}
  {\large\color{faint} #4\par}\vspace{10pt}
  {\small\sffamily\color{faint} #5\par}\vspace{4pt}
  {\color{rule}\rule{\linewidth}{0.4pt}}
  \end{center}\vspace{6pt}
}
% 摘要块: \begin{cnabstract} ... \end{cnabstract}
\newenvironment{cnabstract}{%
  \begin{quote}\small\color{faint}\sffamily\bfseries\color{black} 摘要 \quad\normalfont\rmfamily\color{faint!80!black}%
}{\end{quote}\vspace{2pt}}
% 重点结论框: \begin{quotebox} ... \end{quotebox}
\newenvironment{quotebox}{%
  \begin{quote}\small
  {\color{rule}\rule{\linewidth}{0.8pt}}\par\vspace{2pt}%
}{\par\vspace{2pt}{\color{rule}\rule{\linewidth}{0.8pt}}\end{quote}\vspace{2pt}}

\newcommand{\keywords}[1]{\begin{quote}\small\color{faint}{\bfseries\sffamily\color{black} 关键词:\ } #1\end{quote}}

% 表内彩色标记
\newcommand{\upgood}[1]{\textcolor{good}{\bfseries #1}}
\newcommand{\downbad}[1]{\textcolor{warn}{\bfseries #1}}
\newcommand{\dimtext}[1]{\textcolor{faint}{#1}}

\setlength{\parskip}{2pt}
\linespread{1.12}
\emergencystretch=3em  % CJK 段落容忍度,消除 overfull


\renewcommand{\runningtitle}{LingBot-VLA 2.0 优化全景报告}
\renewcommand{\reportdate}{2026 年 8 月 14 日}
\renewcommand{\reportfooter}{accel-bench · TR-2026-08-14-ALL}

\begin{document}

\reporthead{综合技术报告 · TR-2026-08-14-ALL}{2026 年 8 月 14 日}
{LingBot-VLA 2.0 lerobot 移植与加速:全景报告}
{训练 / 推理最终配置、叠加工程方案与全部实测数字}
{代码:lerobot pr3967-combined \quad 主测平台:8$\times$A100-80GB \quad 辅助:GB10 / 4$\times$4090D(待测)}

\begin{quotebox}
\textbf{一句话定位(2026-08-14 锁定,sdpa 全程保留,从未被否决):}\\
\textbf{训练} = sdpa(自动后端)+ bf16 + 稠密双 GEMM MoE(默认 512 阈值)\newline
\hspace*{1.2em}+ fused AdamW(\texttt{optimizer\_fused=true})+ 不开 grad-ckpt;
单卡 B=3 / DDP 每卡 B=2。\\
\textbf{推理} = sdpa + bf16 + \texttt{compile\_predict\_velocity}\newline
\hspace*{1.2em}(max-autotune-\newline \hspace*{1.2em}no-cudagraphs)+ \texttt{compile\_prefix}
+ \texttt{preprocess\_device=cuda} + \texttt{num\_steps=10},
一次性 bake 进 ckpt 的 config.json,rollout 只需 \texttt{-{}-policy.path}。
\end{quotebox}

\noindent 工作对象:robbyant/lingbot-vla-v2 官方仓与 huggingface/lerobot PR \#3967 移植版
(本地分支 \texttt{pr3967-combined})。全部数字为真实 6B ckpt 实测;主测时宿主机有
LIBERO eval 争抢(load 21--57),同条件前后对比有效,绝对值偏保守,干净机更快。
口径:推理 e2e = 预处理 + model forward(vision+prefix+10 步去噪)+ unapply;
训练 e2e = fwd + bwd + clip + optimizer.step。
推理 workload A = robotwin 3 相机 480$\times$640;B = SO101 单相机。
训练数据集 = gpudad SO101(993 段 / 132k 帧 / 3 相机 512$\times$512)。

\section{推理}

\subsection{叠加工程方案(按落地顺序)}
\begin{table}[htbp]
\centering\footnotesize
\caption{推理优化栈:每一项为什么做、怎么起效、实测结果(A100)}
\begin{tabular}{@{}p{2.9cm}p{3.4cm}p{4.3cm}p{3.6cm}@{}}
\toprule
方案 & 原因 & 原理 & 结果 \\
\midrule
sdpa + bf16;去噪循环不变量外提;MoE 去同步 &
PR 基线 eager,每步去噪重复算不变量、MoE 有隐式同步 &
fused attention kernel;不变量提到循环外;去掉 host 端同步 &
1039.5 $\to$ 951.9 ms(\upgood{-8.4\%}) \\
\addlinespace
\texttt{compile\_predict\_}\newline \texttt{velocity}\newline(max-autotune-\newline no-cudagraphs) &
10 步去噪 = 10 次小图发射,launch 开销主导 &
torch.compile 把去噪步骤融成整图 &
$\to$ 699.5 ms(\upgood{-32.7\%}) \\
\addlinespace
稠密双 GEMM MoE\newline(\texttt{moe\_dense\_}\newline \texttt{max\_tokens=512}) &
小 token 数下 per-expert 循环 kernel 发射多、GEMM 利用率低 &
token 数 $\le$512 时改走两次大 GEMM 的稠密等价计算 &
robotwin 967.4 $\to$ +compile \textbf{398.5 ms}(\upgood{-73.3\%} vs eager 基线);loss 一致到 $9\times10^{-7}$ \\
\addlinespace
\texttt{compile\_prefix}(C3) &
vision+prefix 每次推理重算且未编译 &
prefix KV fill 一并编译(model 口径) &
\upgood{-20\%}(253--269 ms 区间) \\
\addlinespace
\texttt{preprocess\_device}\newline\texttt{=cuda}(C4) &
CPU 预处理占 e2e 近一半(668 ms / 1705 ms) &
图像 resize/normalize 移到 GPU,预处理 \upgood{-62\%} &
e2e \textbf{约 300 ms}(SO101,全程 \upgood{-79\%}) \\
\addlinespace
bake 固化五开关 &
推理开关走 CLI 易漏、训练侧又必须是默认关 &
写进 ckpt \texttt{config.json},\texttt{from\_pretrained} 自动生效 &
rollout 零 flag;一次性编译 43.7 s 后稳态 258 ms \\
\bottomrule
\end{tabular}
\end{table}

保真度(LIBERO drift vs fp32 参考,阈值 2\%):bf16 0.66--0.73\%、dense MoE on/off
0.76\%、C3 0.69\%、C4 0.66\%——全部同量级,无劣化。

\subsection{实验结果}
\begin{table}[htbp]
\centering\small
\caption{推理:官方 upstream vs 本分支(同权重,robotwin 3 相机,10 步去噪)}
\begin{tabular}{@{}lrrr@{}}
\toprule
配置 & e2e (ms) & 预处理 & model \\
\midrule
官方 eager(README 默认) & 1705.0 & 668.1 & 1020.7 \\
官方 use\_compile & 1113.0 & 752.3 & 344.2 \\
\textbf{本分支优化栈} & \textbf{436.6} & 147.8 & \textbf{288.8} \\
\bottomrule
\end{tabular}
\end{table}

e2e 加速 \textbf{3.9$\times$}(对 eager)/ \textbf{2.55$\times$}(对 compile);
model-only 3.53$\times$/1.19$\times$。LIBERO 真实帧路径 293 ms,干净机估算约 250 ms。
官方 README 的 130 ms@4090D claim 在 A100 无法复现(最优 model-only 344 ms),
大概率是 sample\_actions 的 print 计时口径;4$\times$4090D 实测进行中,结果补录。

\begin{table}[htbp]
\centering\small
\caption{推理 attention 消融(锁定栈内,compile+prefix+GPU 预处理,iters=10)}
\begin{tabular}{@{}lrrl@{}}
\toprule
组合 & e2e (ms) & $\Delta$ & 结论 \\
\midrule
\textbf{sdpa + bf16(推荐)} & \textbf{252.9} & \dimtext{—} & 推荐 \\
eager + bf16 & 276.8 & \downbad{+9.4\%} & 慢 \\
eager + fp32(上游原始) & 251.9 & \dimtext{-0.4\%(噪声内)} & compile 下不掉速;训练慢 \\
sdpa + fp32 & 289.7 & \downbad{+14.5\%} & 最差,谁也别用 \\
\bottomrule
\end{tabular}
\end{table}

\section{训练}

\subsection{显存账本(决定一切的前提)}
可训练 5.81B / 冻结 0.42B。固定态 = bf16 参数 11.6GB + 梯度 10.8GB
+ \textbf{AdamW m+v fp32 43.3GB(54\%)}= \textbf{65.7GB};激活约 10GB/样本。
optimizer 状态才是大头,压激活是错杠杆。

\subsection{叠加工程方案}
\begin{table}[htbp]
\centering\footnotesize
\caption{训练优化栈:每一项为什么做、怎么起效、实测结果(A100-80GB)}
\begin{tabular}{@{}p{2.9cm}p{3.4cm}p{4.3cm}p{3.6cm}@{}}
\toprule
方案 & 原因 & 原理 & 结果 \\
\midrule
稠密双 GEMM MoE & 同推理,per-expert 循环在训练 batch 下同样低效 & 双 GEMM 稠密等价 & B=4 fwdbwd 2807.1 $\to$ 1501.3 ms(\upgood{-46.5\%}),loss 一致到 $9\times10^{-7}$ \\
\addlinespace
fused AdamW\newline(\texttt{optimizer\_fused}\newline\texttt{=true}) & 参数逐个 python 循环更新 & 单 kernel 融合 & \upgood{-5.4\%},loss 逐位一致 \\
\addlinespace
sdpa + bf16\newline(不换 eager) & bench 曾显示 eager 快 6.9\% 省 10GB & bench 是相同样本平铺的假象;真实变长序列下 mem-efficient sdpa 更优 & 300 步真实 A/B:sdpa \textbf{0.978 s}/65.7GB vs eager 0.999 s/68.5GB,eager B=4 OOM \\
\addlinespace
不开 grad-ckpt & 省的是不缺钱的激活 & 重算换显存,固定态 65.7GB 纹丝不动 & +38\% 慢,否决 \\
\addlinespace
batch 选择 & 80GB 装得下多少激活 & 65.7GB 固定 + $\sim$10GB/样本 & 单卡 B=3(72.0GB);DDP allreduce +11GB/卡 $\Rightarrow$ 每卡 B=2 \\
\bottomrule
\end{tabular}
\end{table}

\subsection{实验结果}
\begin{table}[htbp]
\centering\small
\caption{单卡训练吞吐排序(e2e = fwd+bwd+clip+step,robotwin 3 相机)}
\begin{tabular}{@{}lrrrl@{}}
\toprule
组合 & step (ms) & samples/s & 峰值显存 & 结论 \\
\midrule
\textbf{B=3 bf16 fused(推荐)} & \textbf{1509.7} & \textbf{1.99} & 72.0GB & \textbf{冠军} \\
B=4 + 8bit AdamW(bnb) & 2159.9 & 1.85 & 73.5GB & 可行但更慢 \\
B=2 fused & 1137.3 & 1.76 & \dimtext{—} & 基线 \\
B=4 + grad-ckpt & \dimtext{—} & \dimtext{更低} & 28.7GB & +38\% 慢,否决 \\
\bottomrule
\end{tabular}
\end{table}

e2e 总收益:PR 原版 B=2 2378.1 ms $\to$ foreach 1202.0 $\to$ \textbf{fused 1137.3 ms}
(\upgood{-52\%},2.09$\times$),三者 loss 0.528622 一致到 $10^{-6}$。

多卡(社区标准 \texttt{accelerate launch} + \texttt{lerobot-train},
\texttt{batch\_size} 为每卡语义):1 卡 1.97 s/step $\to$ 2 卡 1.00 s(4.0 samples/s)
$\to$ 4 卡 1.00 s(\textbf{8.0 samples/s}),2$\to$4 完美线性。
30k 步:单卡 B=3 约 12.6 h,4 卡约 \textbf{8.3 h}。

\begin{table}[htbp]
\centering\small
\caption{训练 attention 长训 A/B(300 步,gpudad 真实数据,2026-08-14)}
\begin{tabular}{@{}lrrr@{}}
\toprule
臂 & updt\_s 稳态 & 峰值显存 & loss 曲线 \\
\midrule
\textbf{LA sdpa+bf16 B=3} & \textbf{0.978} & \textbf{65.7GB} & 完全等价 \\
LB eager+bf16 B=3 & 0.999(\downbad{+2.2\%}) & 68.5GB(\downbad{+2.9GB}) & (max $|\Delta|\approx0.012$) \\
LC eager+bf16 B=4 & OOM & 79.3/80GB & bench 说能跑,真实不行 \\
\bottomrule
\end{tabular}
\end{table}

对官方 upstream 训练:官方单卡 80GB \textbf{完全跑不了}(muon fp32 master+动量
约 74GB,micro=2 + grad-ckpt 仍 OOM,官方设计为 32 卡 fsdp2);官方 4 卡 fsdp2
micro=2/卡(global 8)稳态 2.999 s/step = 0.67 samples/s/GPU;
本分支单卡 2.35 samples/s(fwd+bwd)$\approx$ 官方 per-GPU 吞吐的 \textbf{3.5$\times$}。
口径声明:官方含 MoGe/DINO-video 蒸馏前向支线,移植版没有该开销。

\section{负面结果(全部被实测否决的路)}
\begin{table}[htbp]
\centering\footnotesize
\caption{否决清单:不是 sdpa,是这些东西}
\begin{tabular}{@{}p{4.6cm}p{9.8cm}@{}}
\toprule
被否决项 & 证据 / 原因 \\
\midrule
显式设 \texttt{sdpa\_backend}(任何值) &
CUDNN:bf16+bool mask+GQA 下 backward 1222/1267 梯度 NaN,step 1 即 \texttt{grdn:nan}(前向 parity 完全正常,不足为据);
EFFICIENT:\texttt{No available kernel}\\
eager 训练 & bench 假象(相同样本平铺);300 步真实 A/B:+2.2\% 慢、+2.9GB、B=4 OOM \\
\texttt{attention\_fp32} & +8--14\% 慢,数值无收益 \\
fp16 推理 & launch-bound 无收益(294.7 vs bf16 278.9 ms)且 LIBERO drift 3.33\% 超阈值 \\
flex\_cached 训练 & +75\% 慢 \\
grad-ckpt(B$\le$4) & +38\% 慢,省的是不缺钱的显存 \\
\bottomrule
\end{tabular}
\end{table}

\section{最终结论}
\begin{quotebox}
\textbf{训练}:sdpa(自动后端)+ bf16 + 稠密 MoE + fused AdamW + 不开 grad-ckpt,
单卡 B=3(1.99 samples/s,72GB)/ DDP 每卡 B=2(4 卡 8.0 samples/s,30k 步约 8.3 h)。
相对 PR 原版 \textbf{2.09$\times$},相对官方 per-GPU \textbf{3.5$\times$}。\\
\textbf{推理}:sdpa + bf16 + compile(max-autotune-no-cudagraphs)+ compile\_prefix
+ GPU 预处理 + num\_steps=10,bake 进 ckpt。相对官方 eager \textbf{3.9$\times$}
(1705.0 $\to$ 436.6 ms),SO101 单相机约 300 ms。
\end{quotebox}

配套文档(均在 \texttt{docs/accel/}):
\begin{itemize}[itemsep=0pt,topsep=1pt,parsep=0pt]
\item 操作手册:\texttt{lingbot\_vla\_train\_rollout\_commands.md}
(lerobot-train / lerobot-rollout 命令全文);
\item 开关逐项启动验证(13 臂,11 通过):\texttt{lingbot\_vla\_switch\_table.pdf};
\item reBot B601 实操:\texttt{lingbot\_vla\_rebot\_guide.pdf};
\item 社区英文 quickstart:\newline
\texttt{lingbot\_vla\_v2\_community\_quickstart.md}。
\end{itemize}

\section{已知边界与后续项}
\begin{itemize}[itemsep=0pt,topsep=1pt,parsep=0pt]
\item 全部绝对时延为污染环境(load 21--57)实测,干净机更快;相对结论不受影响。
\item 8 卡 DDP 数字为 2/4 卡线性外推,真 8 卡验证待 GPU 0--3 释放。
\item 训练侧进一步候选:compile 训练整图(-10\~20\%,攻坚)、fsdp2 固定态分片(24GB 卡必需)。
\item 4$\times$4090D 实测(官方 claim 平台)进行中,结果补录。
\end{itemize}

\paragraph{参考与产物} 分支 \texttt{pr3967-combined}(关键 commit:\allowbreak{}\texttt{f2ff2ca5} 稠密 MoE、\texttt{569929b8} C3、\texttt{21ef4187} C4、\texttt{e23a03f5} sdpa\_backend+fused);A100 数据目录 \texttt{\textasciitilde/accel-bench/}(\texttt{train\_e2e/}、\texttt{switch\_smoke/}、\texttt{long\_ab/} 等);备份 \texttt{\textasciitilde/accel-bench/backups/}(git bundle + STATE.md)。

\end{document}
